% #############################################################################
% This is Chapter 5
% !TEX root = ../main.tex
% #############################################################################
\fancychapter{Conclusion}
\label{chap:conclusion}

\section{Summary of the Work}

This dissertation investigated whether Large Language Models can verify the correctness properties of distributed broadcast algorithms, that is, whether a model can decide, from a description of an algorithm and a list of properties, which of those properties are satisfied and which are violated. The problem is relevant because the manual verification of distributed protocols is laborious and error-prone, and because language models have shown a growing ability to reason about code and specifications, so it is natural to ask how far that ability extends to the subtle safety and liveness arguments of broadcast protocols.

The evaluation was conducted iteratively over several phases, each one motivated by the observations of the previous one. The first phase established that among the small local models only Qwen2.5-Coder-7B was both accurate and fully deterministic, which made it a reliable representative for single-run experiments, while other small models were either inaccurate or highly probabilistic and were progressively removed. The second phase showed, through variants that altered only the metadata of an algorithm and not its logic, that the retained model was reasoning about the pseudocode rather than about the name or the declared fault model, although it already displayed some difficulty in separating liveness properties from safety properties. The third phase introduced larger models accessed through APIs and confirmed that they verify correct algorithms at near-perfect rates but do not fundamentally improve the detection of a missing liveness guarantee. The fourth phase introduced genuine logical errors into the pseudocode and revealed that the small and mid-sized models detect essentially none of them, marking every property as satisfied and following the narrative of the description rather than the logic of the protocol. The fifth phase extended the evaluation to higher-potential reasoning models, which detected the injected errors reliably but rejected a large fraction of the correct algorithms, and it also showed, through a confidence analysis, that one of these models is well calibrated while the other is not. Finally, the algorithm set was redesigned into a unified and neutral form with shorter prompts, and this last round both preserved the central findings and produced a small local model that, for the first time, detected the injected errors at a high rate.

\section{Main Findings}

The central finding of this work is an asymmetry between classes of models that defines the current boundary of what language models can do for the verification of distributed algorithms. The small and mid-sized models function as effective confirmation tools, validating correct algorithms with accuracy around ninety to one hundred per cent, but until the final round they detected almost none of the genuine errors. The higher-potential reasoning models are the mirror image, detecting the injected errors with full reliability but rejecting a large fraction of the correct algorithms because they are excessively suspicious. No single model in the study achieved both a high true-positive rate and a low false-negative rate at the same time.

A second finding concerns the properties themselves. The difficulty of verification is not uniform, because the properties that can be checked by inspecting the local state of a process are easier than the properties that require reasoning about the propagation of messages across the whole system and about the behaviour of an adversary. The false positives of the higher-potential model concentrate precisely on Validity, Integrity, and Agreement, which are the properties that demand quantifying over all executions, while a local and syntactic property such as No duplication is verified without error.

A third finding is that the false negatives of the higher-potential model follow a recognizable pattern of incomplete verification, in which the model correctly identifies an attack vector but omits the step of confirming that the quorum or threshold of the protocol prevents the attack from succeeding, or in which the model reasons about a regime that the resilience condition of the algorithm explicitly excludes. On the well calibrated model this incompleteness is accompanied by lower confidence, so that a simple filter that discards low-confidence violations raises the effective accuracy substantially, whereas on the poorly calibrated model the same filter brings no benefit.

A fourth finding, which emerged from the final round, is that the presentation of the algorithm matters. Unifying the pseudocode into a neutral template and shortening the prompts allowed a small local model to detect the injected errors at a rate of ninety per cent that had not been seen in any earlier phase, while at the same time it appears to have removed semantic hints that the higher-potential model had been relying upon on the simplest algorithm. This shows that the way in which a verification task is framed has a measurable effect on the outcome, and that the same change can help one class of models while hurting another.

\section{Limitations}

Several limitations should be kept in mind when interpreting these results. The higher-potential models were accessed through a platform with strict rate limits, which allowed roughly one complete run per day and therefore kept the number of repetitions and the number of valid responses on the broken variants small, so that the detection results for those models, although consistent, rest on a limited sample. In the final round two of the five intended models produced no usable data, one because of an exhausted API balance and the other because its agent had been removed from the platform, which reduced the comparison to three models and left the direct comparison with the higher-potential model incomplete for that round. The set of algorithms and injected errors, while broadened, remains modest in size, and the parser that classifies the responses, although permissive, may misclassify a small number of unusually formatted answers. Finally, the study measures whether a model reaches the correct verdict and gives a consistent justification, and it does not attempt to produce a formal proof or a machine-checkable certificate of correctness.

\section{Future Work}

Several directions follow naturally from this work. The most immediate is to repeat the final round once the higher-potential models are available again, so that the unified algorithm set is evaluated across all model classes under identical conditions, and to enlarge the number of algorithms and injected errors so that the detection results rest on a firmer statistical base. A second direction is to investigate the effect of the algorithm presentation more systematically, by comparing descriptions with and without semantic message labels, since the final round suggested that this choice helps some models and hurts others. A third direction is to exploit the complementary strengths that were observed, by combining the confirmation ability of a small model with the error-detection ability of a higher-potential model and with the confidence signal of the well calibrated model, in order to obtain a verification procedure that is both accurate on correct algorithms and reliable at detecting errors. A final direction is to move from verdicts towards explanations that can be checked, by asking the models to produce the specific execution or the specific quorum argument that supports each verdict, which would make the output auditable and would connect this line of work more closely to formal verification.
